% !PRINTED 257
\weberpage{257}
%<<CONT
Rationalitätsbereich gehören die Funktionen von $\omega$ an, die durch
die beiden Substitutionen

\weq{1}{\left(\omega, -\frac{1}{\omega}\right), \quad (\omega, \omega+2)}

ungeändert bleiben (\S\,54, 2). Wenn also ein System von $\nu$
Funktionen $\Phi_{a,c,\partial}$ durch die Substitutionen (1) nur unter sich
permutiert wird, so sind diese Funktionen die Wurzeln einer
Transformationsgleichung, deren Koeffizienten rational von
$f(\omega)^{24}$ abhängen. Hierzu gehören (wie wir früher schon auf anderem
Wege nachgewiesen haben) für ein \gesperrt{ungerades} $n$, das wir jetzt
immer voraussetzen, gewisse Potenzen der Größen

\[
f\left(\frac{c+\partial\omega}{a}\right),
\]

worin, wie ein- für allemal bemerkt sei, $c$ durch 16, und wenn $n$
nicht durch 3 teilbar ist, durch 48 teilbar angenommen wird.

Die etwas erweiterten Grundsätze des \S\,70 führen verhältnismäßig
 einfach zur Berechnung dieser Gleichungen.

Eine Erweiterung ist aber notwendig aus folgendem Grunde:

Bei den bisherigen Betrachtungen konnten wir den Schluß
machen: wenn eine \gesperrt{rationale} Funktion von $j(\omega)$ für jedes endliche
 $\omega$ mit positiv imaginärem Teil endlich bleibt, so ist sie
eine \gesperrt{ganze} Funktion von $j(\omega)$, weil zu jedem \gesperrt{endlichen} $j(\omega)$
auch ein endliches, nicht reelles $\omega$ gehört (\S\,52). Bei den rationalen
 Funktionen von $f(\omega)$ können wir aber aus der Endlichkeit
für jedes endliche imaginäre $\omega$ nur schließen, daß sie ganze
Funktionen von $f(\omega)$ und $1:f(\omega)$ sind, weil nur zu jedem endlichen
 $f(\omega)$ mit Ausnahme von $f(\omega) = 0$ ein endliches imaginäres
$\omega$ gehört.

Es entspricht aber der Substitution

\weq{2}{\left(\omega, \frac{\omega-1}{\omega+1}\right)}

die Vertauschung

\[
\left(f(\omega), \frac{\sqrt{2}}{f(\omega)}\right), \qquad [\S\,34, (18)]
\]

und wenn also die Funktionen $\Phi_{a,c,\partial}$ so gewählt sind, daß sie
auch durch die Substitution (2) nur untereinander permutiert
werden, so werden die Koeffizienten der Gleichung, deren Wurzeln
sie sind, rational von
%CONT>>
